% GENERATED FILE. Edit canonical lessons, then run npm run generate:books.
% canonical-source-hash: fnv1a64:6680a7ab297f05d3
% canonical-lessons: FR-C299-chute, FR-C299-fouler, FR-C299-contusion, FR-C299-hematome, FR-C299-plaie, FR-R299-revision-292-299

\chapter{Une mauvaise chute}
\label{ch:fr-une-mauvaise-chute}
\hlchaptermodality{299}

\begin{chapteropening}
\textbf{By the end of this chapter:} \emph{I can describe a fall and what it does: a fall, spraining something, a bruise, a large bruise and a wound.}

Complete the last lesson of chapter 299: I can describe a fall and what it does: a fall, spraining something, a bruise, a large bruise and a wound.
\end{chapteropening}

\section[la chute]{la chute — the fall; the drop}

\label{lesson:FR-C299-chute}



\begin{quote}
Before the new one: say the French for the airbag, then the French for to dent.
\end{quote}

\subsection*{You'll want to know: la chute}
\textbf{la chute} — \textquotedblleft{}the fall; the drop\textquotedblright{}.

\begin{quote}
Ma grand-mère a fait une chute dans l'escalier.
\end{quote}

\begin{quote}
\emph{My grandmother had a fall on the stairs.}
\end{quote}

\emph{Les chutes de neige} — snowfalls; \emph{une chute d'eau} — a waterfall.

\begin{cousinweb}
\emph{chute} is built on \emph{chu}, the old past participle of \emph{choir} \textquotedblleft{}to fall\textquotedblright{} (EtymDB), which is inherited from Latin \emph{cadere} \textquotedblleft{}to fall\textquotedblright{} (Lewis \& Short).
\end{cousinweb}

\begin{grammarlens}[title={What you've built}]
One more word for telling, explaining and describing at B1.
\end{grammarlens}

\subsection*{Your turn}
\begin{itemize}
\raggedright
  \item \emph{Say it:} \emph{la chute}
  \item \emph{Say it:} \emph{la chute}, once more
  \item \emph{Say it:} a sentence of your own with \emph{la chute}
  \item \emph{Recall:} say the French for the airbag, then the French for to dent, then say \emph{la chute} again
\end{itemize}

\begin{tcolorbox}[breakable,colback=teal!4,colframe=teal!35!black,title={Before you move on}]
What does la chute mean? (\textquotedblleft{}The fall; the drop\textquotedblright{}.) Say it once more.
\end{tcolorbox}
\section[fouler]{fouler — to sprain (as se fouler); to tread on}

\label{lesson:FR-C299-fouler}



\begin{quote}
Before the new one: say the French for to dent, then the French for the fall.
\end{quote}

\subsection*{You'll want to know: fouler}
\textbf{fouler} — \textquotedblleft{}to sprain (as se fouler); to tread on\textquotedblright{}.

\begin{quote}
Il s'est foulé la cheville en courant.
\end{quote}

\begin{quote}
\emph{He sprained his ankle while running.}
\end{quote}

In this sense mostly reflexive, with a body part: \emph{se fouler le poignet}. Without \emph{se}, it means to tread on or to press: \emph{fouler le sol}.

\begin{cousinweb}
\emph{fouler} is inherited through Old French \emph{foler} \textquotedblleft{}to crush, to trample\textquotedblright{} from a verb built on Latin \emph{fullō} \textquotedblleft{}a fuller\textquotedblright{}, who trod cloth (EtymDB).
\end{cousinweb}

\begin{grammarlens}[title={What you've built}]
One more word for telling, explaining and describing at B1.
\end{grammarlens}

\subsection*{Your turn}
\begin{itemize}
\raggedright
  \item \emph{Say it:} \emph{fouler}
  \item \emph{Say it:} \emph{fouler}, once more
  \item \emph{Say it:} a sentence of your own with \emph{fouler}
  \item \emph{Recall:} say the French for to dent, then the French for the fall, then say \emph{fouler} again
\end{itemize}

\begin{tcolorbox}[breakable,colback=teal!4,colframe=teal!35!black,title={Before you move on}]
What does fouler mean? (\textquotedblleft{}To sprain (as se fouler); to tread on\textquotedblright{}.) Say it once more.
\end{tcolorbox}
\section[la contusion]{la contusion — the bruise, the contusion}

\label{lesson:FR-C299-contusion}



\begin{quote}
Before the new one: say the French for the fall, then the French for to sprain.
\end{quote}

\subsection*{You'll want to know: la contusion}
\textbf{la contusion} — \textquotedblleft{}the bruise, the contusion\textquotedblright{}.

\begin{quote}
Il a seulement quelques contusions.
\end{quote}

\begin{quote}
\emph{He only has a few bruises.}
\end{quote}

Medical; in everyday French, \emph{un bleu}.

\begin{cousinweb}
\emph{contusion} is a learned borrowing from Latin \emph{contūsiō} \textquotedblleft{}a bruising\textquotedblright{} (EtymDB), from \emph{contundere} \textquotedblleft{}to beat, to bruise\textquotedblright{}.
\end{cousinweb}

\begin{grammarlens}[title={What you've built}]
One more word for telling, explaining and describing at B1.
\end{grammarlens}

\subsection*{Your turn}
\begin{itemize}
\raggedright
  \item \emph{Say it:} \emph{la contusion}
  \item \emph{Say it:} \emph{la contusion}, once more
  \item \emph{Say it:} a sentence of your own with \emph{la contusion}
  \item \emph{Recall:} say the French for the fall, then the French for to sprain, then say \emph{la contusion} again
\end{itemize}

\begin{tcolorbox}[breakable,colback=teal!4,colframe=teal!35!black,title={Before you move on}]
What does la contusion mean? (\textquotedblleft{}The bruise, the contusion\textquotedblright{}.) Say it once more.
\end{tcolorbox}
\section[l'hématome]{l'hématome — the haematoma, the large bruise}

\label{lesson:FR-C299-hematome}



\begin{quote}
Before the new one: say the French for to sprain, then the French for the bruise.
\end{quote}

\subsection*{You'll want to know: l'hématome}
\textbf{l'hématome} — \textquotedblleft{}the haematoma, the large bruise\textquotedblright{}.

\begin{quote}
Elle a un gros hématome sur la jambe.
\end{quote}

\begin{quote}
\emph{She has a big bruise on her leg.}
\end{quote}

Masculine: \emph{un hématome}.

\begin{cousinweb}
\emph{hématome} is a learned medical word built on Greek \emph{haîma} \textquotedblleft{}blood\textquotedblright{} with the suffix \emph{-ōma} (EtymDB lists no parent).
\end{cousinweb}

\begin{grammarlens}[title={What you've built}]
One more word for telling, explaining and describing at B1.
\end{grammarlens}

\subsection*{Your turn}
\begin{itemize}
\raggedright
  \item \emph{Say it:} \emph{l'hématome}
  \item \emph{Say it:} \emph{l'hématome}, once more
  \item \emph{Say it:} a sentence of your own with \emph{l'hématome}
  \item \emph{Recall:} say the French for to sprain, then the French for the bruise, then say \emph{l'hématome} again
\end{itemize}

\begin{tcolorbox}[breakable,colback=teal!4,colframe=teal!35!black,title={Before you move on}]
What does l'hématome mean? (\textquotedblleft{}The haematoma, the large bruise\textquotedblright{}.) Say it once more.
\end{tcolorbox}
\section[la plaie]{la plaie — the wound, the cut}

\label{lesson:FR-C299-plaie}



\begin{quote}
Before the new one: say the French for the bruise, then the French for the haematoma.
\end{quote}

\subsection*{You'll want to know: la plaie}
\textbf{la plaie} — \textquotedblleft{}the wound, the cut\textquotedblright{}.

\begin{quote}
L'infirmière a nettoyé la plaie.
\end{quote}

\begin{quote}
\emph{The nurse cleaned the wound.}
\end{quote}

Informal: \emph{quelle plaie !} — what a pain!

\begin{cousinweb}
\emph{plaie} is inherited from Latin \emph{plāga} \textquotedblleft{}a blow, a wound\textquotedblright{} (EtymDB).
\end{cousinweb}

\begin{grammarlens}[title={What you've built}]
One more word for telling, explaining and describing at B1.
\end{grammarlens}

\subsection*{Your turn}
\begin{itemize}
\raggedright
  \item \emph{Say it:} \emph{la plaie}
  \item \emph{Say it:} \emph{la plaie}, once more
  \item \emph{Say it:} a sentence of your own with \emph{la plaie}
  \item \emph{Recall:} say the French for the bruise, then the French for the haematoma, then say \emph{la plaie} again
\end{itemize}

\begin{tcolorbox}[breakable,colback=teal!4,colframe=teal!35!black,title={Before you move on}]
What does la plaie mean? (\textquotedblleft{}The wound, the cut\textquotedblright{}.) Say it once more.
\end{tcolorbox}
\section[(dialogue)]{Retour sur les chapitres 292–299}

\label{lesson:FR-R299-revision-292-299}



\begin{quote}
Say the French for the investigator and the inspector.
\end{quote}

\subsection*{Your turn}
Say these aloud:

\begin{itemize}
\raggedright
  \item the investigator, the inspector, the private detective, the clue, the print
  \item the mystery, to solve, to prove, the alibi, the suspicion
  \item to stop and question, the arrest, the handcuffs, police custody, to question
  \item to confess, the confession, the accomplice, to acquit, to release
  \item the armed robbery, the armed robber, the hostage, the ransom, to escape
  \item the pile-up, the reckless driver, speeding, the accident report, the skid
  \item to knock over, the injured person, the shock, the airbag, to dent
  \item the fall, to sprain, the bruise, the haematoma, the wound
\end{itemize}

\begin{tcolorbox}[breakable,colback=teal!4,colframe=teal!35!black,title={Before you move on}]
The investigator? (\textbf{l'enquêteur}.) The inspector? (\textbf{l'inspecteur}.) The private detective? (\textbf{le détective}.) The clue? (\textbf{l'indice}.) The print? (\textbf{l'empreinte}.) The mystery? (\textbf{le mystère}.) To solve? (\textbf{élucider}.) To prove? (\textbf{prouver}.) The alibi? (\textbf{l'alibi}.) The suspicion? (\textbf{le soupçon}.) To stop and question? (\textbf{interpeller}.) The arrest? (\textbf{l'arrestation}.) The handcuffs? (\textbf{les menottes}.) Police custody? (\textbf{la garde à vue}.) To question? (\textbf{interroger}.) To confess? (\textbf{avouer}.) The confession? (\textbf{l'aveu}.) The accomplice? (\textbf{le complice}.) To acquit? (\textbf{acquitter}.) To release? (\textbf{relâcher}.) The armed robbery? (\textbf{le braquage}.) The armed robber? (\textbf{le braqueur}.) The hostage? (\textbf{l'otage}.) The ransom? (\textbf{la rançon}.) To escape? (\textbf{échapper}.) The pile-up? (\textbf{le carambolage}.) The reckless driver? (\textbf{le chauffard}.) Speeding? (\textbf{l'excès de vitesse}.) The accident report? (\textbf{le constat}.) The skid? (\textbf{le dérapage}.) To knock over? (\textbf{renverser}.) The injured person? (\textbf{le blessé}.) The shock? (\textbf{le choc}.) The airbag? (\textbf{l'airbag}.) To dent? (\textbf{cabosser}.) The fall? (\textbf{la chute}.) To sprain? (\textbf{fouler}.) The bruise? (\textbf{la contusion}.) The haematoma? (\textbf{l'hématome}.) The wound? (\textbf{la plaie}.)
\end{tcolorbox}
